\documentclass[12pt, twoside]{article}
%\documentclass[12pt, twoside]{article}
\usepackage[letterpaper, margin=1in, headsep=0.2in]{geometry}
\setlength{\headheight}{0.6in}
%\usepackage[english]{babel}
\usepackage[utf8]{inputenc}
\usepackage{microtype}
\usepackage{amsmath}
\usepackage{amssymb}
%\usepackage{amsfonts}
\usepackage[nomessages]{fp} %\FPeval{\var-name}{2*sin(pi/6)}
\usepackage{siunitx} %units in math. eg 20\milli\meter
\usepackage{yhmath} % for arcs, overparenth command
\usepackage{tikz} %graphics
\usetikzlibrary{quotes, angles, arrows, arrows.meta}
\usepackage{pgfplots}
\usepgfplotslibrary{fillbetween}
\pgfplotsset{compat=1.16}
\usepackage{graphicx} %consider setting \graphicspath{{images/}}
\usepackage{parskip} %no paragraph indent
\usepackage{enumitem}
\usepackage{multicol}
\usepackage{venndiagram}

\usepackage{fancyhdr}
\pagestyle{fancy}
\fancyhf{}
\renewcommand{\headrulewidth}{0pt} % disable the underline of the header
\raggedbottom
\hfuzz=2mm %suppresses overfull box warnings

\usepackage{hyperref}

\fancyhead[LE]{\thepage}
\fancyhead[RO]{Markscheme}
\fancyhead[LO]{La Scuola d'Italia / Dr.\ Huson / IB Math 12 \\* 13 October 2026}

\begin{document}
\subsubsection*{1.10 Pre-Quiz: Statistics --- Markscheme}

\begin{enumerate}[itemsep=0.3cm]

%--- Problem 1: Grouped frequency table; continuous data; modal class; estimated mean and sd; stratified sample [8 marks] ---
\item[1.] [Maximum mark: 8]

Mid-interval values $1, 3, 5, 7, 9$ with frequencies $6, 14, 20, 12, 8$;
$n = 60$.

\emph{\textbf{GDC:} enter the mid-interval values in one list and the
frequencies in a second list, then run \texttt{One-Variable Statistics} with
the frequency list.}

\textbf{(a)} Continuous \hfill \textbf{A1}

\emph{Time is measured, not counted. Accept a reason such as ``time can
take any value in an interval''; no reason is required.}

\textbf{(b)} $4 \le h < 6$ \hfill \textbf{A1}

\emph{The class, not its frequency: ``$20$'' alone earns A0. Accept
``$4$ to $6$ hours''. The mid-interval value $5$ alone earns A0.}

\textbf{(c)} $\bar{x} = \dfrac{6(1) + 14(3) + 20(5) + 12(7) + 8(9)}{60}$
\hfill \textbf{(M1)}

$= \dfrac{304}{60} = 5.0666\ldots = 5.07$ hours (3 s.f.) \hfill \textbf{A1}

\emph{M1 for using the mid-interval values with the frequencies, shown as a
sum or as the GDC lists. Accept $5.067$ or $\frac{76}{15}$. Using the lower
class boundaries ($0, 2, 4, 6, 8$, giving $4.07$) or upper boundaries
(giving $6.07$) earns M0A0. The mean of the frequencies, $12$, earns no
marks.}

\textbf{(d)} The exact number of hours for each student is not known; each
student is taken to be at the mid-interval value of their class.
\hfill \textbf{R1}

\emph{R1 for either idea: the raw data are not given (only the classes), or
mid-interval values were used in place of the actual values.}

\textbf{(e)} $\sigma = 2.3371\ldots = 2.34$ hours (3 s.f.) \hfill \textbf{A1}

\emph{Do not accept $s_{x} = 2.3568\ldots$ ($2.36$), the sample standard
deviation.}

\textbf{(f)} $\dfrac{112}{480} \times 60$ \hfill \textbf{M1}

$= 14$ students \hfill \textbf{A1}

\emph{M1 for the Grade 12 proportion of the school applied to $60$, however
written (e.g.\ $\frac{60}{480} = \frac{1}{8}$, then $\frac{112}{8}$).}

\newpage

%--- Problem 2: Five-number summary from raw data; IQR; outlier test; draw the box plot [9 marks] ---
\item[2.] [Maximum mark: 9]

Ordered: $21, 22, 23, 24, 25, 25, 26, \mathbf{26}, 27, 27, 28, 29, 30, 31,
44$ \quad ($n = 15$).

\emph{\textbf{GDC:} the TI-Nspire excludes the median from each half, so
$Q_{1}$ is the median of the lowest seven values and $Q_{3}$ the median of
the highest seven.}

\textbf{(a)} Median $= 26$ \hfill \textbf{A1}

$Q_{1} = 24$, \quad $Q_{3} = 29$ \hfill \textbf{A1}

\emph{Both quartiles are needed for the second A1. Accept $Q_{1} = 24.5$ and
$Q_{3} = 28.5$ from a hand method that includes the median in each half;
feedback: the GDC convention is the one to use.}

\textbf{(b)} $\text{IQR} = 29 - 24 = 5$ \hfill \textbf{A1}

\emph{FT: award A1 for $Q_{3} - Q_{1}$ from the candidate's own quartiles.
The range, $44 - 21 = 23$, earns no mark.}

\textbf{(c)} $Q_{3} + 1.5 \times \text{IQR} = 29 + 1.5 \times 5$
\hfill \textbf{M1}

$= 36.5$ \hfill \textbf{A1}

$44 > 36.5$, so $44$ is an outlier. \hfill \textbf{R1}

\emph{M1 for $1.5 \times$ their IQR added to their $Q_{3}$. R1 for the
comparison with $44$ stated, not just the fence value. Adding
$1.5 \times \text{IQR}$ to the median or the mean earns no marks. The lower
fence, $24 - 7.5 = 16.5$, is not required. FT: award all three marks from the
candidate's own $Q_{3}$ and IQR, provided their fence is below $44$.}

\textbf{(d)}

\begin{center}
\begin{tikzpicture}
  \begin{scope}[x=0.35cm]
    \draw (15,0)--(51,0);
    \foreach \t in {15,16,...,51} \draw (\t,-0.08)--(\t,0.08);
    \foreach \t in {15,20,...,50} \node[below] at (\t,-0.12) {\footnotesize \t};
    \draw[thick] (24,0.6) rectangle (29,1.4);
    \draw[thick] (26,0.6)--(26,1.4);
    \draw[thick] (21,1.0)--(24,1.0);
    \draw[thick] (29,1.0)--(31,1.0);
    \draw[thick] (21,0.8)--(21,1.2);
    \draw[thick] (31,0.8)--(31,1.2);
    \node at (44,1.0) {$\times$};
  \end{scope}
\end{tikzpicture}
\end{center}

Box from $Q_{1} = 24$ to $Q_{3} = 29$ with the median line at $26$
\hfill \textbf{A1}

Whiskers to $21$ and $31$ \hfill \textbf{A1}

Outlier $44$ marked with a cross \hfill \textbf{A1}

\emph{Award each mark independently; plotted values within half a small
division ($\pm 0.5$ minute) are accepted. A whisker drawn to $44$ or to
$36.5$ loses the second A1. FT: award the first A1 for a box drawn at the
candidate's own quartiles from (a).}

\newpage

%--- Problem 3: Read median, IQR, a count and a percentile from a CF curve; frequencies from the curve [9 marks] ---
\item[3.] [Maximum mark: 9]

$80$ students; the curve passes through $(20, 8)$, $(30, 20)$, $(40, 40)$,
$(55, 60)$ and $(60, 65)$.

\emph{Readings within one small square of the values below are accepted
(\,$\pm 2.5$ seconds on the time axis, $\pm 1$ student on the cumulative
frequency axis), and any value following from the candidate's own readings
earns the A1. Feedback: draw the reading lines on the graph.}

\textbf{(a)} Median: cumulative frequency $40$ gives $t = 40$ seconds
\hfill \textbf{A1}

\textbf{(b)} $Q_{1} = 30$ (at $20$), \quad $Q_{3} = 55$ (at $60$)
\hfill \textbf{(M1)}

$\text{IQR} = 55 - 30 = 25$ seconds \hfill \textbf{A1}

\emph{M1 for reading the times at cumulative frequencies $20$ and $60$ (or
for either quartile correct). Reading at $25$ and $75$ (percent of $100$)
earns M0.}

\textbf{(c)} $80 - 65$ \hfill \textbf{(M1)}

$= 15$ students \hfill \textbf{A1}

\emph{M1 for reading the cumulative frequency at $t = 60$ and subtracting
from $80$. The reading $65$ alone, as a final answer, earns M1A0.}

\textbf{(d)} $10\% \text{ of } 80 = 8$ \hfill \textbf{(M1)}

Cumulative frequency $8$ gives $t = 20$ seconds \hfill \textbf{A1}

\emph{M1 for $8$ (or for reading the time at cumulative frequency $8$).
Reading at cumulative frequency $72$ (the slowest $10\%$, $t = 70$) earns
M1A0.}

\textbf{(e)} $p = 40 - 8 = 32$ \hfill \textbf{A1}

$q = 65 - 40 = 25$ \hfill \textbf{A1}

\emph{The two marks are independent. Accept $q = 80 - 8 - 32 - 15 = 25$.
FT: award the A1 for $q$ from the candidate's own $p$ when $q$ is found by
subtraction from $80$.}

\newpage

%--- Problem 4: Pearson's r, describing correlation, y-on-x regression line, interpreting the gradient, interpolation, extrapolation [9 marks] ---
\item[4.] [Maximum mark: 9]

$n = 8$, \quad $\bar{x} = 7$, \quad $\bar{y} = 68$.

\emph{\textbf{GDC:} enter $x$ and $y$ in two lists and run
\texttt{Linear Regression (mx+b)}.}

\textbf{(a)} $r = 0.98815\ldots = 0.988$ (3 s.f.) \hfill \textbf{A1}

\emph{Do not accept $r^{2} = 0.976$, which sits beside $r$ on the
regression screen.}

\textbf{(b)} Strong, positive (linear) correlation \hfill \textbf{A1A1}

\emph{A1 for ``strong'' (accept ``very strong''), A1 for ``positive''; the two
marks are independent. ``Moderate'' or ``weak'' earns A0 for the strength.
FT from the candidate's $r$ in (a).}

\textbf{(c)} $y = 3.72x + 42.0$ \hfill \textbf{A1A1}

\emph{A1 for $a = 3.72$, A1 for $b = 42.0$; the two marks are independent.
Accept $a = 3.7173\ldots$, $b = 41.978\ldots$ (or $42$), or either value to
more figures. The reversed line $x = 0.263y - 10.9$ earns A0A0.}

\textbf{(d)} For each extra hour of revision, the test score increases by
about $3.72$ points. \hfill \textbf{A1}

\emph{The answer must give the rate (about $3.7$ points) per hour, in
context. ``Students who revise more score higher'' restates the correlation
without the rate and earns A0. FT from the candidate's $a$ in (c).}

\textbf{(e)} $y = 3.72(10) + 42.0$ \hfill \textbf{(M1)}

$= 79.152\ldots \approx 79.2$ \hfill \textbf{A1}

\emph{Accept any value consistent with the candidate's own correctly-rounded
coefficients, including $79.152\ldots$ (full precision) and $79.2$ (from
$3.72$ and $42.0$), and either rounded to $79$. FT from the candidate's
line in (c). \textbf{Feedback:} store the coefficients in the GDC (or use the
regression function \texttt{f1(10)}) and round only the final answer.}

\textbf{(f)} $20$ hours is outside the range of the data ($2$ to $12$ hours),
so the estimate would be an extrapolation and is not reliable.
\hfill \textbf{R1}

\emph{R1 for stating that $20$ lies outside the data range (or ``above the
greatest number of hours recorded''). Accept the observation that the line
gives $116$, which is impossible for a score out of $100$, as a supporting
reason. ``Revising too much might not help'' with no reference to the data
earns R0.}

\end{enumerate}
\end{document}
